% =====================================================================================================
% Online Appendix F: Economic Implications, additional results (supports Section V). Writer: economics writer (v3),
% 2026-09-24; round-3 revision by WP5, 2026-09-25 (fix list E2, E5; the fleet share moved here from Section V.C).
% Module rb3_econ2 (memo + review; round-3 changes in output/memos/rb5_econ.md). Tables:
% paper/tables/appG_econ_{growth,wall,gamma}.tex (copies of output/tables/rb3_econ2_app_*.tex); figure:
% output/figures/rb3_econ2_figure.pdf. Input in main.tex after appendix_wedge (Online Appendix F).
% =====================================================================================================

\section[Economic Implications: Additional Results]{Appendix \thesection. Economic Implications: Additional Results}\label{app:econ}

This appendix supports Section~\ref{sec:econ}. Every estimate is regenerated by one deterministic script in the
replication package from the clean sample and decision units of Section~\ref{sec:wedge}, the technologies and IsoFLOP
minima of Section~\ref{sec:tech}, and two vintages of the Epoch AI database; values taken from the literature are cited
where they appear.

\textit{Compute growth and the vintage of the database.} Table~\ref{tab:app-econ-growth} applies one
estimator, the running top ten of training compute at release with a wild cluster bootstrap by developer, to Epoch AI's
public file as captured on May 31, 2024 \citep{epochai2024dataarchive} and to the snapshot of September 23, 2026
\citep{epochai2026data}. On the 2024 file the estimator gives 4.13-fold growth, consistent with the published 4.2-fold
rate of \citet{sevilla2024training} ($p=0.73$). A two-factor Shapley decomposition attributes 64 percent of the gap to
the 2026 file's 5.23-fold rate to revised compute estimates and 36 percent to entries added after May 2024.

\textit{The trend in the revealed wedge.} Table~\ref{tab:app-econ-trend}, panel~A, gives the trend under alternative
estimators: within developers, above $10^{24}$ FLOP, weighted by compute, for the four disclosed flagship runs above
$10^{25}$ FLOP, and on version 3's 56 family-label units. Over the decision units every interval includes no growth,
except that of the within-developer trend conditional on log compute (2.70-fold a year). Panel~B gives the
compute-weighted wedge by release year, and panel~C the five-year multiples of the scenario of Table~\ref{tab:econ},
panel~B, when the frontier's wedge keeps rising.

\textit{The data wall at observed allocations.} A developer with wedge $w$ minimizes lifetime cost $6ND+\Xi N$, where
$\Xi N$ is its value of compactness, subject to reaching the loss of its unconstrained choice; $\Xi$ is set so that the
unconstrained choice has $\varepsilon_N/\varepsilon_D=w$. Effective data follow \citet{muennighoff2023scaling}, and
the technology is the reference or a member of its observationally equivalent family with the same expansion path and
frontier. All objects depend on $w$ and scarcity only. Table~\ref{tab:app-econ-wall} reports the extra lifetime cost by
$w$ and $r$ (panel~A), the measured wedge under a cap (panel~B), today's frontier under each technology and at the stock's
lower bound (panel~C), and, for compute-optimal runs, the hard wall, the shadow value, the elasticity of substitution
between compute and unique tokens, and the return to compute behind the wall (panel~D). Figure~\ref{fig:app-econ} plots
the main relations. Two sensitivities are worth stating. If the frontier's wedge had followed the open-weight trend since
the median release date of the four anchor runs (mid-2024), frontier runs would already process 3.5 to 8.9 times the
compute-optimal tokens ($\sigma^*=0.60$ to 0.74; 6.2 at 0.70). They would have reached the 100-trillion-token stock
between late 2025 and early 2026, so starting the rise in September 2026, as Section~\ref{sec:econ} does, is the
conservative choice. And in the full repetition model the decay of excess parameters acts as a size penalty, so that
measured wedges of models near the expansion path rise rather than fall as the wall binds (for the frontier with its
wedge held, $\hat w/w$ is 1.04 in 2029 and 1.11 in 2030, against 0.94 and 0.87 under data decay only). Because data are
nonrival \citep{jones2020nonrivalry}, the private shadow values of unique tokens in Section~\ref{sec:econ-wall}
understate their social value when several developers use the same tokens.

\textit{The fleet share of AI compute.} Read as planned serving expenditure, an interpretation that the serving evidence
of Section~\ref{sec:wedge} does not support, the 2025 compute-weighted wedge of the decision units ($w-1=4.4$) implies a
share of AI compute spent on inference. Suppose that each model is served for one year while training compute grows
5.06-fold a year, that a FLOP costs the same in serving as in training, and that final training runs are 10 to 23
percent of developers' research compute \citep{denain2026final}. The implied inference share is then 0.17 to 0.33 (0.68
for one-year lives if all AI compute were final training runs). Longer service lives lower it, because the models in
service are then older vintages trained with less compute: it is 0.11 to 0.23 with two-year lives and 0.08 to 0.17 with
three-year lives (Table~\ref{tab:app-econ-gamma}, panel~C). Without Alibaba the one-year range is 0.09 to 0.19. Disclosures are higher: inference took about
three-fifths of Google's machine-learning energy in 2019--2021 \citep{patterson2022carbon} and 70 percent of Meta's AI
power capacity \citep{wu2022sustainable}, and it was about 40 percent of Nvidia's data-center revenue in fiscal 2024
\citep{nvidia2024q4call}. The comparison is a loose check, and the 2025 wedge rests largely on one family, Alibaba's
Qwen3.

\textit{Growth calibration, fleet shares and synthetic tokens.} Table~\ref{tab:app-econ-gamma} reports $\gamma$ and
the compute needed to halve reducible loss for every large design (panel~A); profiles of the IsoFLOP frontier in the
irreducible loss $E$ (panel~B), which bound $E$ tightly for Chinchilla and Llama~3 but loosely or not at all for Marin's
two-decade ladders, and of which only Chinchilla's is in known loss units and contains its design's joint-fit $E$; the
fleet share of inference implied by the 2025 wedge for service lives of one to three years, R\&D multiples of final
training runs and relative prices of an inference FLOP (panel~C); and the break-even scarcities at which generating a
token costs its shadow value (panel~D).

\input{tables/appG_econ_growth}

% Online Appendix table (round 3), copied from output/tables/rb3_econ2_app_trend.tex (rb3_econ2) by WP5, 2026-09-25;
% regenerate with code/analysis/rb3_econ2/run.py and re-copy (the copy adds only these two comment lines).
% Online Appendix table -- the trend in the revealed wedge and the uncapped frontier scenario.
% Module rb3_econ2; sources rb3_econ2_wedge_trend.csv, _wedge_trend_lodo.csv, _wedge_trend_by_tech.csv,
% _wedge_by_year.csv, _scenarios_by_trend.csv.
\begin{table}[tp]
\centering
\caption{The Trend in the Revealed Wedge and Frontier Data Demand}
\label{tab:app-econ-trend}
\footnotesize
\setlength{\tabcolsep}{2pt}
\begin{tabular*}{\textwidth}{@{\extracolsep{\fill}}lrrccc@{}}
\toprule
\multicolumn{6}{@{}l}{\textit{Panel A. Growth per year of the revealed wedge and the data multiple, 2023--2025}} \\[2pt]
 & & & & \multicolumn{2}{c}{$D/D^*$ per year} \\
 & & & & at $\sigma^*=0.70$ & at $\sigma^*=$ \\
\cmidrule(l){5-6}
Estimator & $n$ & $w$ per year & 95\% CI & (from $M/M^*$) & 0.60; 0.74 \\
\midrule
Decision units, OLS (primary) & 49 & 1.86 & [0.96, 3.54] & 2.06 & 1.59; 2.41 \\
\quad developer fixed effects & 49 & 1.59 & [0.47, 5.59] & 1.72 & 1.42; 1.94 \\
\quad fixed effects and log compute & 49 & 2.70 & [1.39, 5.21] & 3.19 & 2.10; 4.10 \\
\quad above $10^{24}$ FLOP (6 developers) & 11 & 1.76 & [0.78, 4.15] & 1.93 & 1.53; 2.23 \\
\quad compute-weighted & 49 & 2.21 & [0.38, 13.1] & 2.53 & 1.81; 3.09 \\
Flagship runs above $10^{25}$ FLOP$^a$ & 4 & 1.70 & -- & 1.86 & 1.49; 2.12 \\
Models, OLS & 77 & 1.99 & [0.95, 4.33] & 2.24 & 1.68; 2.67 \\
Models, compute-weighted & 77 & 1.96 & [1.01, 3.85] & 2.19 & 1.66; 2.60 \\
Aggregate wedge, decision units & 49 & 2.02 & -- & 2.28 & 1.69; 2.72 \\
Median wedge, decision units & 49 & 1.53 & -- & 1.64 & 1.37; 1.83 \\
Family-label units (version 3), OLS & 56 & 1.91 & [1.03, 3.48] & 2.13 & 1.62; 2.51 \\
Leave one developer out (units, OLS) & -- & 1.66--2.26 & -- & -- & -- \\
32 technologies (models, OLS) & 77 & 1.17--4.28 & -- & 1.97--2.38 & -- \\
\midrule
\multicolumn{6}{@{}l}{\textit{Panel B. The compute-weighted wedge by release year}} \\[2pt]
 & & & & \multicolumn{2}{c}{$D/D^*$ at this $w$} \\
\cmidrule(l){5-6}
Release year & Units & $s$ & $w$ & $\sigma^*=0.70$ & 0.60; 0.74 \\
\midrule
2023 & 16 & 0.24 & 1.32 & 1.38 & 1.23; 1.48 \\
2024 & 27 & 0.54 & 2.17 & 2.46 & 1.79; 3.00 \\
2025 & 6 & 0.81 & 5.38 & 7.13 & 3.53; 10.98 \\
\midrule
\multicolumn{6}{@{}l}{\textit{Panel C. The frontier adopts an open-weight trend without a stop}} \\[2pt]
 & & $D/D^*$ & \multicolumn{2}{c}{$D$ over 5 years} & Years to \\
\cmidrule(lr){4-5}
Trend & $\sigma^*$ & per year & $4.2\times$ & $5.06\times$ & 2025 $w$ \\
\midrule
Decision units & 0.60 & 1.59 & 359 & 571 & 2.4 \\
 & 0.70 & 2.06 & 1,305 & 2,074 & 2.4 \\
 & 0.74 & 2.41 & 2,887 & 4,588 & 2.4 \\
Within developers & 0.70 & 1.72 & 529 & 841 & 3.2 \\
Units above $10^{24}$ FLOP & 0.70 & 1.93 & 945 & 1,502 & 2.6 \\
Compute-weighted units & 0.70 & 2.52 & 3,580 & 5,690 & 1.9 \\
Flagships above $10^{25}$ FLOP$^a$ & 0.70 & 1.86 & 774 & 1,230 & 2.8 \\
\bottomrule
\end{tabular*}
\vspace{2pt}
\begin{tablenotes}[Notes]
Panel A: growth factor $\exp(\hat b)$ of the slope of $\ln w$ (reference technology) on release date; decision units as in Section~\ref{sec:wedge}; 95 percent wild cluster bootstrap-$t$ intervals by developer (Webb weights, $B=9{,}999$; 18 developers; 6 above $10^{24}$ FLOP, so that interval is indicative); developer fixed effects: the slope is identified by the 9 developers with units at more than one date; $^a$the four disclosed dense runs above $10^{25}$ FLOP in 2024--2025, descriptive (OLS, no interval; 1.31 without Pangu Ultra); aggregate and median wedge: growth between the 2023 and 2025 values, with the aggregate $w=1/(1-s)$ at the compute-weighted $s$; $D/D^*$ from $M/M^*$: $\exp(\hat b/2\theta)$, $\theta=1/\sigma^*-1$ of the reference, which needs no curvature given $M^*(C)$ and equals the conversion at the reference $\sigma^*=0.70$; $\sigma^*=0.60/0.74$: $w$ growth converted with $D/D^*=w^{\sigma^*/[2(1-\sigma^*)]}$; 32 technologies: the ex-ante technologies of Section~\ref{sec:wedge}, each in its own parameter convention, on the audited token counts (the reference gives the models row). Panel B: $s=\sum(w^+-1)C/\sum w^+C$ over decision units, $w^+=\max(w,1)$, reference technology; $w=1/(1-s)$; $D/D^*$ that a model with this wedge would use at given compute; over models, $s$ is 0.27, 0.58, 0.83. Panel C: the scenario of Table~\ref{tab:econ}, panel~B, without the stop at the 2025 wedge; $D$ over 5 years: frontier tokens in 2031 relative to September 2026; years until the frontier's wedge reaches the 2025 compute-weighted wedge of the decision units ($w=5.38$).
\end{tablenotes}
\end{table}


\input{tables/appG_econ_wall}

\input{tables/appG_econ_gamma}

\begin{figure}[tp]
\centering
\includegraphics[width=\textwidth]{rb3_econ2_figure}
\caption{Data Demand with Rising Wedges, the Data Wall at Observed Allocations, and the Shadow Value of Data}
\label{fig:app-econ}
\begin{figurenotes}
Reference technology ($\kappa$ free, $\sigma^*=0.70$) unless stated. Panel a: training tokens of frontier runs at
5.06-fold compute growth, anchored at the median data multiple of the four disclosed dense runs above $10^{25}$ FLOP;
black, wedge held; colored, wedge rising 1.86-fold a year from September 2026 (the trend over the decision units)
with $D/D^*=w^{\sigma^*/[2(1-\sigma^*)]}$; dotted, the rise stops at the 2025 aggregate wedge ($w=5.38$). Gray dots: dense models in the verified sample above
$10^{24}$ FLOP. Solid and dashed gray lines: the quality-adjusted (100T) and repetition-adjusted (320T) stocks of
public text of \citet{villalobos2022run}, growing 5 percent a year; shading: the 95 percent interval of the unique stock
with stock growth between 0 and 10 percent a year. Panel b: extra lifetime cost against compute-optimal scarcity
$r=D^*(C)/U$ for $w=1$, 2 and 5.38; solid, data decay only ($R^*_D=15.4$); dashed, data-only fit ($R^*_D=2.9$). Panel c:
the wedge measured from the capped choice, relative to one plus the value of compactness at that choice (solid) and to
the no-wall wedge (dashed), for $w=5.38$, against the no-wall tokens over the stock; the hard-cap curve is the token-cap
case of Proposition~\ref{prop:wedge}. Panel d: shadow value of a unique token for a compute-optimal run at the frontier
compute of 2026 (at \$$10^{-18}$ per FLOP, amortized, and at cloud prices) and at $10^{26}$ FLOP; shading: posted API
prices of DeepSeek-V4.1-Flash (off-peak) to DeepSeek-V4-Pro (peak) per million output tokens on September 24, 2026
\citep{deepseek2026pricing}.
\end{figurenotes}
\end{figure}
